\documentclass[10pt,a4paper]{amsart}
\usepackage[margin=19mm]{geometry}
\usepackage{amsmath,amssymb,mathtools}
\usepackage[scr=rsfs,cal=euler]{mathalfa}
\usepackage{xfrac}
\usepackage[unicode,hidelinks,bookmarks=false]{hyperref}
\newcommand{\ie}{i.e.\ }
\newcommand{\cf}{cf.\ }
\newcommand{\resp}{respectively\ }
\newcommand{\Subsection}{\subsection}
\providecommand{\nobreakdash}{\nobreak\hbox{-}}
\setlength{\emergencystretch}{2em}
\multlinegap0pt
\usepackage{bm}
\usepackage{bbm}
\usepackage{dsfont}
\usepackage{xspace}
\usepackage{cancel}
\usepackage{mathtools}
\usepackage{amsthm}
\makeatletter
\@ifundefined{theorem}{\newtheorem{theorem}{Theorem}}{}
\@ifundefined{proposition}{\newtheorem{proposition}[theorem]{Proposition}}{}
\makeatother
\title[On Extending Type $B$ Parking Spaces]{On Extending Type $B$ Parking Spaces}
\author{Anthony Adams, Joshua Dorsam, Lily Levitsky, Megan Mann}
\date{Source: arXiv:2601.18090v1 (CC BY 4.0)}
\begin{document}
\maketitle

\noindent\emph{Source and licence.} On Extending Type $B$ Parking Spaces. Version arXiv:2601.18090v1 (CC BY 4.0), distributed under the Creative Commons Attribution 4.0 International licence, as recorded in the arXiv licence metadata for this exact version and shown on the abstract page https://arxiv.org/abs/2601.18090v1. Full rights notice: licenses/arxiv-2601.18090-CC-BY-4.0.txt.\par\smallskip

\noindent\textit{Editorial scope.} Counted unit: the Theorem whose source label is \texttt{thm:OPn1Thm} in the article (source0.tex lines 387--391, source label \texttt{thm:OPn1Thm}), delivered together with the article's own complete original proof of it (source0.tex lines 393--438, one \texttt{proof} environment of 46 lines). No mathematics is written, repaired or re-derived: statement and proof are the article's own text line for line. Required context retained uncounted: OPdecomp (lines 321--328). Every definition, display and label of those lines is retained unchanged. Omitted from this package and not redistributed: 1--320, 329--386, 392--392, 439--631. Nothing inside a retained excerpt is omitted or altered: every retained body is delivered exactly as the article's lines are written, so no omission receipt of any kind accompanies this package. Citations: the retained excerpts carry no citation command at all, so no bibliography entry of the article is transcribed and no reference is redistributed. Reference adaptation: none was needed and none was made; every reference inside the delivered unit resolves inside it, so no unresolved reference or citation is produced. Printed slips kept literal: no printed word, symbol, hypothesis or number of the counted statement or of its proof was altered. Layout and class: the article's own class is \texttt{amsart} with packages including graphicx, amssymb, xcolor, ulem, amsthm, biblatex, hyperref; the wrapper is the lane's pinned \texttt{amsart} editorial wrapper at 10pt on A4, which restates the theorem-like environments the retained text needs and adds the article's own macro definitions that the retained text needs (none), with extra wrapper packages (bm, bbm, dsfont, xspace, cancel, mathtools). The delivered numbering is the unit's own. Assets: no retained span contains a figure environment, a graphics-inclusion command, a PDF-page inclusion, a table or a TikZ picture; no image file is copied into this package, the package declares no asset dependency and no image file is redistributed with it. Licence: version arXiv:2601.18090v1 of this article is distributed under the Creative Commons Attribution 4.0 International licence, as recorded in the arXiv licence metadata for this exact version and shown on the abstract page https://arxiv.org/abs/2601.18090v1. A full rights notice is delivered at licenses/arxiv-2601.18090-CC-BY-4.0.txt. Characters: every retained excerpt is pure ASCII, so the delivered engine, XeLaTeX with its default Latin Modern text font, reports no missing character.
\begin{proposition}\label{OPdecomp}
   $$
        OP_{n, k} \cong \bigoplus_{(\lambda,\mu) \vDash n} d(\lambda, k+1) d(\mu, k) V_{\lambda,\mu}
   $$
    where the sum is over bipartitions $(\lambda,\mu)$ of $n$ where $\lambda$ has at most $k+1$ rows and $\mu$ has at most $k$ rows, and the Weyl dimension formula gives
    $$d(\lambda,a)=\begin{cases} \frac{\prod_{i < j \leq a} (\lambda_i - \lambda_j+j-i)}{\prod_{i<j \leq a}{(j-i)}}. & \text{rows}(\lambda)\le a \\
    0 & \text{rows}(\lambda)>a\end{cases}$$
\end{proposition}

\begin{theorem} \label{thm:OPn1Thm}
For all $n,$ the $B_{n+1}-$representation
$$\bigoplus_{(\lambda,a)\in X_{n,m}}  (\lceil \frac{\lambda_1 + 1 - \lambda_2}{3} \rceil ) V_{(\lambda',a)}$$
restricts to $OP_{n, 1}.$ In fact, it is the only $B_{n+1}-$extension of $OP_{n,1}$ which decomposes into irreducibles indexed by $\tilde{X}_{n,3}.$ 
\end{theorem}

\begin{proof}
Recall by proposition \ref{OPdecomp} that 
$$  
         OP_{n,1}\cong \bigoplus_{i=0}^n \bigoplus_{\lambda \vdash i} (\lambda_1 - \lambda_2 + 1)V_{((\lambda_1, \lambda_2),n-i)}, 
   $$

so for  $k=1 $ ($m=3$), we have that 
$$X_{n,3}=\{ ((\lambda_1,\lambda_2),a)\vDash n \}$$ 
and 
$$\tilde X_{n,3}= \{(\lambda',a)| (\lambda,a)\in X_{n,3}\}.$$

 Thus, finding an extension to $B_{n+1}$ comprised of irreducible representations corresponding to $\tilde X_{n,3}$ is equivalent to finding coefficients  $c_{\lambda' }$ satisfying the following equation:
\begin{equation}\label{eq: extension problem for OPn1}
    \bigoplus_{i=0}^n \bigoplus_{\lambda \vdash i} (\lambda_1 - \lambda_2 + 1)V_{((\lambda_1, \lambda_2),n-i)} \cong \bigoplus_{(\lambda, a) \in X_n} c_{\lambda'}\mathrm{Res}_{B_n}^{B_{n+1}} V_{\lambda',a}.
\end{equation}   
We will find a recursive formula for $c_{\lambda'}$ in terms of $\lambda_1-\lambda_2,$ as only $\lambda_1-\lambda_2$ determines the coefficient of $V_{\lambda,a}$ in the Weyl dimension formula. 


First suppose $\lambda_1=\lambda_2.$ Then the only irreducible representation corresponding to  $\tilde{X}_{n,3}$ whose restriction contains $V_{\lambda,a}$ is $V_{\lambda',a},$ since partitions  $(\mu,b)\in \tilde X_{n,1}$ have $(\mu_1>\mu_2)$ by construction. Hence, $$c_{\lambda'}= \lambda_1-\lambda_2+1=1.$$

Suppose $\lambda_1-\lambda_2=1.$ There are exactly two bipartitions in $\tilde X_{n,3}$ such that removing one box yields $(\lambda,a): (\lambda',a)$ and $(\lambda,a+1).$ Hence, 
\begin{align*}
    2 &= \lambda_1-\lambda_2+1\\
    &= c_{\lambda}+c_{\lambda'}.
\end{align*}
But recall that $\lambda_1-1=\lambda_2,$ so by the previous case,
$$ c_{\lambda}= c_{(\lambda_1-1,\lambda_2)'}=1.$$ Therefore, 
$$c_{\lambda'}=1.$$ 

Suppose $\lambda_1-\lambda_2\ge 2.$ Then a box can be added to either row of $\lambda$ to obtain a partition in $\tilde X_{n,3}.$ Hence there are three bipartitions in $\tilde X_{n,3}$ such that removing one box yields $(\lambda,a):$ $(\lambda',a), ((\lambda_1,\lambda_2+1),a),$ and $(\lambda,a+1).$ This gives the recursive formula  

\begin{equation} 
\lambda_1-\lambda_2 +1 = c_{\lambda'}+ c_{(\lambda_1-1,\lambda_2)'} + c_{(\lambda_1-1,\lambda_2+1)'} .
\label{eq:recursion}
\end{equation}
 It follows by induction that if $\lambda_1-\lambda_2=\sigma_1-\sigma_2$, then $c_{\lambda'}=c_{\sigma'}.$ Therefore, for any nonnegative integer $a$ we can define $c_a=c_{\lambda'}$ for any partition $\lambda$ with $\lambda_1-\lambda_2+1=a. $ For $\lambda$ with $\lambda_1-\lambda_2\ge 2$ transforms the recursion above into the form 

 $$c_a= \begin{cases} 1 & a\le 2\\
    a- c_{a-1}-c_{a-2} &\text{rows}(\lambda)>2
\end{cases}$$
This recurrence relation has a unique solution $c_a=\lceil \frac{a}{3}\rceil.$  Therefore, we indeed have the desired decomposition

\begin{equation*} OP_{n,1}= \bigoplus_{(\lambda,a)\in X_{n,m}}  (\lceil \frac{\lambda_1 + 1 - \lambda_2}{3} \rceil ) V_{(\lambda',a)}.
\qedhere
\end{equation*}
\end{proof}

\end{document}
